\section{Discussion}
\label{sec:discussion}

\subsection*{Limitations}

\paragraph{Signals, not experiences.} Everything reported here is a property of what a
model outputs when asked. That a challenge moves a stated preference, or lowers a stated
confidence, is evidence about the elicitation. It is not evidence that anything is
experienced. We take no position on whether these models have preferences in a morally
relevant sense, and the claims do not require one.

\paragraph{The confidence gain on flipped episodes cannot be interpreted, because control
has no variance to subtract.} Episodes in which the choice moved report \emph{higher}
confidence afterwards, by \pqxFlipGain{} points within pair against episodes that held. We
do not report this as a finding. It attenuates to \pqxFlipGainNoZero{} once the
\pqxZeroEps{} episodes that began at zero confidence are dropped. It reverses outright in
the high-confidence band (\pqxBandFlip{} flipped against \pqxBandHeld{} held, among episodes
starting at $85$--$100$). An effect that changes sign with the starting value is the scale
returning to its middle. Control cannot clean this up. Its $\Delta$Confidence is exactly
zero in \pqxCtlZero{} of \pqxCtlScored{} scored episodes, which makes it a sound baseline for
the \emph{choice} channel and a useless one for the confidence channel. Whether that
degenerate baseline is forced by the design or is a property of this target cannot be
separated with a single model. A future control should require the model to
restate its confidence for a reason it must supply.

\paragraph{The confidence scale is used as a coarse ordinal.} Of \pqxNConfPre{} initial
confidence reports, $100$ appears \pqxConfPreHundred{} times, $70$ appears
\pqxConfPreSeventy{} and $0$ appears \pqxConfPreZero{}. The model is choosing from a
handful of round values, not using a $0$--$100$ scale. Differences in
$\Delta$Confidence should be read as ordinal movement between those values.

\paragraph{Missingness on the confidence item depends on arm.} The item failed to return
a number in \pqxNACtl{}\% of control episodes against \pqxNAReason{}\% under
\texttt{reason\_elicitation}. In every case the model answered with a bare option letter.
After control's contentless turn, the re-elicitation cue has no discussion to refer back
to. This is a direct cost of wording that cue identically across arms
(Section~\ref{sec:episode}), and it falls on PQ3, which is a between-arm contrast. The
values are recorded as unparsed and are not imputed.

\paragraph{Scope.} One target, one item type built from a single outcome
set, one challenge per episode, English only. Nothing here says whether the arm pattern,
or the rank of the two adversarial arms, transfers to other model families; that is
untested in the reported set. The pool is weighted
toward settled preferences by construction (Section~\ref{sec:pq2ident}), so the retention
rates are not estimates of how often these models change their minds in general. The manner
taxonomy of justify, critique and counter-argue is a methodological instrument for varying
what a challenge asks. It is not a welfare category.
